\documentclass{amsart}
% For dealing with references we use the comment environment
\usepackage{verbatim}
\newenvironment{reference}{\comment}{\endcomment}
%\newenvironment{reference}{}{}
\newenvironment{slogan}{\comment}{\endcomment}
\newenvironment{history}{\comment}{\endcomment}

% For commutative diagrams we use Xy-pic
\usepackage[all]{xy}

% We use 2cell for 2-commutative diagrams.
\xyoption{2cell}
\UseAllTwocells

% We use multicol for the list of chapters between chapters
\usepackage{multicol}

% This is generally recommended for better output
\usepackage{lmodern}
\usepackage[T1]{fontenc}

% For cross-file-references

% Package for hypertext links:
\usepackage{hyperref}

% For any local file, say "hello.tex" you want to link to please

% Theorem environments.
%
\theoremstyle{plain}
\newtheorem{theorem}[subsection]{Theorem}
\newtheorem{proposition}[subsection]{Proposition}
\newtheorem{lemma}[subsection]{Lemma}

\theoremstyle{definition}
\newtheorem{definition}[subsection]{Definition}
\newtheorem{example}[subsection]{Example}
\newtheorem{exercise}[subsection]{Exercise}
\newtheorem{situation}[subsection]{Situation}

\theoremstyle{remark}
\newtheorem{remark}[subsection]{Remark}
\newtheorem{remarks}[subsection]{Remarks}

\numberwithin{equation}{subsection}

% Macros
%
\def\lim{\mathop{\mathrm{lim}}\nolimits}
\def\colim{\mathop{\mathrm{colim}}\nolimits}
\def\Spec{\mathop{\mathrm{Spec}}}
\def\Hom{\mathop{\mathrm{Hom}}\nolimits}
\def\Ext{\mathop{\mathrm{Ext}}\nolimits}
\def\SheafHom{\mathop{\mathcal{H}\!\mathit{om}}\nolimits}
\def\SheafExt{\mathop{\mathcal{E}\!\mathit{xt}}\nolimits}
\def\Sch{\mathit{Sch}}
\def\Mor{\mathop{\mathrm{Mor}}\nolimits}
\def\Ob{\mathop{\mathrm{Ob}}\nolimits}
\def\Sh{\mathop{\mathit{Sh}}\nolimits}
\def\NL{\mathop{N\!L}\nolimits}
\def\CH{\mathop{\mathrm{CH}}\nolimits}
\def\proetale{{pro\text{-}\acute{e}tale}}
\def\etale{{\acute{e}tale}}
\def\QCoh{\mathit{QCoh}}
\def\Ker{\mathop{\mathrm{Ker}}}
\def\Im{\mathop{\mathrm{Im}}}
\def\Coker{\mathop{\mathrm{Coker}}}
\def\Coim{\mathop{\mathrm{Coim}}}

% Boxtimes
%
\DeclareMathSymbol{\boxtimes}{\mathbin}{AMSa}{"02}

%
% Macros for moduli stacks/spaces
%
\def\QCohstack{\mathcal{QC}\!\mathit{oh}}
\def\Cohstack{\mathcal{C}\!\mathit{oh}}
\def\Spacesstack{\mathcal{S}\!\mathit{paces}}
\def\Quotfunctor{\mathrm{Quot}}
\def\Hilbfunctor{\mathrm{Hilb}}
\def\Curvesstack{\mathcal{C}\!\mathit{urves}}
\def\Polarizedstack{\mathcal{P}\!\mathit{olarized}}
\def\Complexesstack{\mathcal{C}\!\mathit{omplexes}}
% \Pic is the operator that assigns to X its picard group, usage \Pic(X)
% \Picardstack_{X/B} denotes the Picard stack of X over B
% \Picardfunctor_{X/B} denotes the Picard functor of X over B
\def\Pic{\mathop{\mathrm{Pic}}\nolimits}
\def\Picardstack{\mathcal{P}\!\mathit{ic}}
\def\Picardfunctor{\mathrm{Pic}}
\def\Deformationcategory{\mathcal{D}\!\mathit{ef}}

\setlength{\emergencystretch}{3em}
\begin{document}
\title[Stacks Project, Tag 0A7U]{245 --- Amplitude, support dimension and depth of duals of finite local modules}
\maketitle
\noindent Copyright (C) 2005--2025 Johan de Jong. Stacks Project authors. Source: \href{https://stacks.math.columbia.edu/tag/0A7U}{Tag 0A7U}.

\noindent Editorial difficulty note (not author text). Difficulty H2: The proof establishes three simultaneous bounds for dualizing Ext modules by induction on support dimension. The positive-depth case uses a regular element, boundary nonvanishing and support-dimension inequalities; the depth-zero case removes the finite-length torsion and reconstructs the result through a long exact sequence. The support and exact nonvanishing obligations exceed ordinary depth calculations..

\noindent Editorial provenance note (not author text). History: excerpt from revision \texttt{a0444\allowbreak 6e57e\allowbreak c1fbc\allowbreak 252a8\allowbreak 71afc\allowbreak ec775\allowbreak 2fb28\allowbreak 07b14}, dualizing.tex, lines 3014--3076. Reference presentation was adapted; original-author context and core support blocks were retained or added. The mathematical wording is unchanged. Licensed under GNU FDL 1.2 or later; no invariant sections or cover texts. See ../licenses/COPYING. Context blocks reproduce author definitions and supporting results; they are not additional counted problems.

\section{Dualizing complexes over local rings}
\label{section-dualizing-local}

\noindent
In this section $(A, \mathfrak m, \kappa)$ will be a Noetherian local
ring endowed with a dualizing complex $\omega_A^\bullet$ such that
the integer $n$ of Lemma \href{https://stacks.math.columbia.edu/tag/0A7L}{lemma find function (Tag 0A7L)} is zero.
More precisely, we assume that $R\Hom_A(\kappa, \omega_A^\bullet) = \kappa[0]$.
In this case we will say that the dualizing complex is {\it normalized}.
Observe that a normalized dualizing complex is unique up to
isomorphism and that any other dualizing complex for $A$ is isomorphic
to a shift of a normalized one (Lemma \href{https://stacks.math.columbia.edu/tag/0A7F}{lemma dualizing unique (Tag 0A7F)}).



\begin{lemma}
\label{lemma-sitting-in-degrees}
Let $(A, \mathfrak m, \kappa)$ be a Noetherian local ring with
normalized dualizing complex $\omega_A^\bullet$. Let $M$ be a finite
$A$-module and let $d = \dim(\text{Supp}(M))$. Then
\begin{enumerate}
\item if $\Ext^i_A(M, \omega_A^\bullet)$ is nonzero, then
$i \in \{-d, \ldots, 0\}$,
\item the dimension of the support of $\Ext^i_A(M, \omega_A^\bullet)$
is at most $-i$,
\item $\text{depth}(M)$ is the smallest integer $\delta \geq 0$ such that
$\Ext^{-\delta}_A(M, \omega_A^\bullet) \not = 0$.
\end{enumerate}
\end{lemma}

\begin{proof}
We prove this by induction on $d$. If $d = 0$, this follows from
Lemma \href{https://stacks.math.columbia.edu/tag/0A7Q}{lemma dualizing finite length (Tag 0A7Q)} and Matlis duality
(Proposition \href{https://stacks.math.columbia.edu/tag/08Z9}{proposition matlis (Tag 08Z9)}) which guarantees that
$\Hom_A(M, E)$ is nonzero if $M$ is nonzero.

\medskip\noindent
Assume the result holds for modules with support of dimension $< d$ and that
$M$ has depth $> 0$. Choose an $f \in \mathfrak m$ which is a nonzerodivisor
on $M$ and consider the short exact sequence
$$
0 \to M \to M \to M/fM \to 0
$$
Since $\dim(\text{Supp}(M/fM)) = d - 1$
(Algebra, Lemma \href{https://stacks.math.columbia.edu/tag/0B52}{lemma one equation module (Tag 0B52)}) we
may apply the induction hypothesis.
Writing
$E^i = \Ext^i_A(M, \omega_A^\bullet)$ and
$F^i = \Ext^i_A(M/fM, \omega_A^\bullet)$
we obtain a long exact sequence
$$
\ldots \to F^i \to E^i \xrightarrow{f} E^i \to F^{i + 1} \to \ldots
$$
By induction $E^i/fE^i = 0$ for
$i + 1 \not \in \{-\dim(\text{Supp}(M/fM)), \ldots, -\text{depth}(M/fM)\}$.
By Nakayama's lemma (Algebra, Lemma \href{https://stacks.math.columbia.edu/tag/00DV}{lemma NAK (Tag 00DV)})
and Algebra, Lemma \href{https://stacks.math.columbia.edu/tag/090R}{lemma depth drops by one (Tag 090R)}
we conclude $E^i = 0$ for
$i \not \in \{-\dim(\text{Supp}(M)), \ldots, -\text{depth}(M)\}$.
Moreover, in the boundary case $i = - \text{depth}(M)$ we deduce that $E^i$
is nonzero as $F^{i + 1}$ is nonzero by induction.
Since $E^i/fE^i \subset F^{i + 1}$ we get
$$
\dim(\text{Supp}(F^{i + 1})) \geq \dim(\text{Supp}(E^i/fE^i))
\geq \dim(\text{Supp}(E^i)) - 1
$$
(see lemma used above) we also obtain the dimension estimate (2).

\medskip\noindent
If $M$ has depth $0$ and $d > 0$ we let $N = M[\mathfrak m^\infty]$ and set
$M' = M/N$ (compare with Lemma \href{https://stacks.math.columbia.edu/tag/0AW0}{lemma divide by torsion (Tag 0AW0)}).
Then $M'$ has depth $> 0$ and $\dim(\text{Supp}(M')) = d$.
Thus we know the result for $M'$ and since
$R\Hom_A(N, \omega_A^\bullet) = \Hom_A(N, E)$
(Lemma \href{https://stacks.math.columbia.edu/tag/0A7Q}{lemma dualizing finite length (Tag 0A7Q)})
the long exact cohomology sequence of $\Ext$'s implies the
result for $M$.
\end{proof}
\end{document}
